\documentclass[12pt,letterpaper]{article}
\usepackage[utf8]{inputenc}
\usepackage{amsmath,amssymb}
\usepackage{helvet}
\renewcommand{\familydefault}{\sfdefault}
\usepackage[margin=1.2in]{geometry}

\title{\textbf{Technical Communication in Cybersecurity}}
\author{Paarth Pandey \\ MNIT Jaipur}
\date{\today}

\begin{document}
\maketitle

\section{Overview}
Effective documentation is critical for IT/OT infrastructure planning and threat mitigation. \LaTeX{} provides a deterministic framework for generating immutable security reports and compliance documentation.

\section{Cryptographic Typesetting}
Technical documentation often requires formatting asymmetric encryption algorithms. For instance, the RSA encryption standard is mathematically represented as:
\begin{equation}
C \equiv M^e \pmod{n}
\end{equation}

Inline representations, such as evaluating computational complexity $O(n \log n)$ or defining hash lengths $H(x) \in \{0, 1\}^{256}$, maintain high typographical precision.

\section{Documentation Benefits}
\begin{itemize}
    \item Version-controlled source files via Git.
    \item Standardized rendering across varied environments.
    \item Robust bibliography management for threat intelligence citations.
\end{itemize}
\end{document}
